\documentclass[tikz,border=3mm,multi=tikzpicture]{standalone}
\usepackage{eleve-geometria}

\begin{document}

% FIGURA 1 — fig-01-translacao-por-vetor
\begin{tikzpicture}[eleve figura, x=1cm, y=1cm]
  \path[use as bounding box] (-0.20,-0.20) rectangle (10.20,5.55);
  \coordinate (A) at (0.75,0.75);
  \coordinate (B) at (3.30,0.75);
  \coordinate (C) at (1.45,3.20);
  \coordinate (Ap) at (6.20,1.60);
  \coordinate (Bp) at (8.75,1.60);
  \coordinate (Cp) at (6.90,4.05);

  \draw[eleve preenchimento, line width=1.5pt] (A) -- (B) -- (C) -- cycle;
  \draw[eleve destaque, fill=orange!14, line width=1.5pt]
    (Ap) -- (Bp) -- (Cp) -- cycle;
  \EleveVertice{A}{below left}{A}
  \EleveVertice{B}{below right}{B}
  \EleveVertice{C}{above}{C}
  \fill[elevelaranja] (Ap) circle (2pt) node[font=\large, text=elevelaranja, below left] {$A'$};
  \fill[elevelaranja] (Bp) circle (2pt) node[font=\large, text=elevelaranja, below right] {$B'$};
  \fill[elevelaranja] (Cp) circle (2pt) node[font=\large, text=elevelaranja, above] {$C'$};

  \foreach \X/\Y in {A/Ap,B/Bp,C/Cp}{
    \draw[eleve auxiliar, ->, line width=1.15pt] (\X) -- (\Y);
  }
  \draw[eleve destaque, ->, line width=2.2pt] (3.75,4.65) -- (6.20,5.02)
    node[midway, above=4pt, font=\Large\bfseries, text=elevelaranja] {$\vec v$};
\end{tikzpicture}

% FIGURA 2 — fig-02-reflexao-por-eixo
\begin{tikzpicture}[eleve figura, x=1cm, y=1cm]
  \path[use as bounding box] (-0.20,-0.15) rectangle (9.20,6.35);
  \draw[eleve destaque, <->, line width=1.8pt] (4.50,0.25) -- (4.50,5.90)
    node[above, font=\Large\bfseries, text=elevelaranja] {$e$};

  \coordinate (A) at (1.05,1.05);
  \coordinate (B) at (3.05,1.05);
  \coordinate (C) at (2.15,4.35);
  \coordinate (Ap) at (7.95,1.05);
  \coordinate (Bp) at (5.95,1.05);
  \coordinate (Cp) at (6.85,4.35);
  \draw[eleve preenchimento, line width=1.5pt] (A)--(B)--(C)--cycle;
  \draw[eleve destaque, fill=orange!14, line width=1.5pt] (Ap)--(Bp)--(Cp)--cycle;

  \foreach \L/\R in {A/Ap,B/Bp,C/Cp}{
    \draw[eleve auxiliar, line width=1.05pt] (\L) -- (\R);
  }
  \draw[elevecinza, line width=1pt] (4.50,1.05) -- (4.78,1.05) -- (4.78,1.33);
  \draw[elevecinza, line width=1pt] (4.50,4.35) -- (4.78,4.35) -- (4.78,4.63);

  \EleveVertice{A}{below left}{A}
  \EleveVertice{B}{below right}{B}
  \EleveVertice{C}{above}{C}
  \fill[elevelaranja] (Ap) circle (2pt) node[font=\large, text=elevelaranja, below right] {$A'$};
  \fill[elevelaranja] (Bp) circle (2pt) node[font=\large, text=elevelaranja, below left] {$B'$};
  \fill[elevelaranja] (Cp) circle (2pt) node[font=\large, text=elevelaranja, above] {$C'$};
  \node[font=\large\bfseries, text=elevecinza] at (4.50,0.02) {mesma distância};
\end{tikzpicture}

% FIGURA 3 — fig-03-rotacao-por-centro
\begin{tikzpicture}[eleve figura, x=1cm, y=1cm]
  \path[use as bounding box] (-0.25,-0.25) rectangle (8.25,7.45);
  \coordinate (O) at (3.95,2.55);
  \coordinate (P) at ($(O)+(15:2.75)$);
  \coordinate (Pp) at ($(O)+(105:2.75)$);
  \coordinate (Q) at ($(O)+(15:1.45)$);
  \coordinate (Qp) at ($(O)+(105:1.45)$);
  \coordinate (R) at ($(O)+(48:2.05)$);
  \coordinate (Rp) at ($(O)+(138:2.05)$);

  \draw[eleve auxiliar] (O) circle (2.75);
  \draw[eleve auxiliar] (O)--(P);
  \draw[eleve auxiliar] (O)--(Pp);
  \draw[eleve preenchimento, line width=1.5pt] (P)--(Q)--(R)--cycle;
  \draw[eleve destaque, fill=orange!14, line width=1.5pt] (Pp)--(Qp)--(Rp)--cycle;
  \fill[eleveazul] (O) circle (2.4pt) node[font=\large\bfseries, text=elevecinza, below] {$O$};
  \fill[eleveazul] (P) circle (2pt) node[font=\large, text=eleveazul, right] {$P$};
  \fill[elevelaranja] (Pp) circle (2pt) node[font=\large, text=elevelaranja, above] {$P'$};

  \draw[eleve destaque, ->, line width=1.8pt]
    ($(O)+(15:1.05)$) arc[start angle=15,end angle=105,radius=1.05];
  \node[font=\Large\bfseries, text=elevelaranja] at (4.46,3.58) {$\alpha$};
  \node[font=\large\bfseries, text=elevecinza] at (3.95,6.78) {$OP=OP'$};
\end{tikzpicture}

% FIGURA 4 — fig-04-ordem-da-composicao
\begin{tikzpicture}[eleve figura, x=1cm, y=1cm]
  \path[use as bounding box] (-0.20,-0.20) rectangle (9.20,8.20);
  \node[font=\large\bfseries, text=elevecinza] at (4.50,7.78)
    {eixo em $0$ \quad deslocamento $+4$};

  % Primeira ordem: 1 -> -1 -> 3.
  \node[font=\Large\bfseries, text=eleveazul] at (4.50,6.85)
    {refletir $\rightarrow$ transladar};
  \draw[eleve auxiliar, ->, line width=1.4pt] (1.55,5.35) -- (4.05,5.35);
  \draw[eleve auxiliar, ->, line width=1.4pt] (4.55,5.35) -- (7.05,5.35);
  \fill[eleveazul] (1.30,5.35) circle (4pt);
  \fill[elevecinza] (4.30,5.35) circle (4pt);
  \fill[elevelaranja] (7.30,5.35) circle (4pt);
  \node[font=\Large\bfseries, text=eleveazul] at (1.30,5.90) {$P=1$};
  \node[font=\Large\bfseries, text=elevecinza] at (4.30,4.80) {$-1$};
  \node[font=\Large\bfseries, text=elevelaranja] at (7.30,5.90) {final $=3$};

  % Segunda ordem: 1 -> 5 -> -5.
  \node[font=\Large\bfseries, text=eleveazul] at (4.50,3.55)
    {transladar $\rightarrow$ refletir};
  \draw[eleve auxiliar, ->, line width=1.4pt] (1.55,2.05) -- (4.05,2.05);
  \draw[eleve auxiliar, ->, line width=1.4pt] (4.55,2.05) -- (7.05,2.05);
  \fill[eleveazul] (1.30,2.05) circle (4pt);
  \fill[elevecinza] (4.30,2.05) circle (4pt);
  \fill[elevelaranja] (7.30,2.05) circle (4pt);
  \node[font=\Large\bfseries, text=eleveazul] at (1.30,2.60) {$P=1$};
  \node[font=\Large\bfseries, text=elevecinza] at (4.30,1.50) {$5$};
  \node[font=\Large\bfseries, text=elevelaranja] at (7.30,2.60) {final $=-5$};
\end{tikzpicture}

% FIGURA 5 — fig-05-tesselacao-hexagonal
\begin{tikzpicture}[eleve figura, x=1cm, y=1cm]
  \path[use as bounding box] (-0.20,-0.20) rectangle (8.20,7.65);
  \coordinate (V) at (4.00,3.55);

  % Três hexágonos regulares compartilham o vértice V.
  \foreach \ang/\cor in {30/eleveazulclaro,150/orange!16,270/eleveazulclaro}{
    \coordinate (C) at ($(V)+(\ang:1.73)$);
    \draw[eleve aresta, fill=\cor, line width=1.4pt]
      ($(C)+(\ang+180:1.73)$)
      \foreach \k in {1,...,5} { -- ($(C)+(\ang+180+60*\k:1.73)$) } -- cycle;
  }
  \fill[elevelaranja] (V) circle (3pt);

  \draw[eleve destaque, line width=1.6pt]
    ($(V)+(0:0.72)$) arc[start angle=0,end angle=120,radius=0.72];
  \draw[eleve destaque, line width=1.6pt]
    ($(V)+(120:0.72)$) arc[start angle=120,end angle=240,radius=0.72];
  \draw[eleve destaque, line width=1.6pt]
    ($(V)+(240:0.72)$) arc[start angle=240,end angle=360,radius=0.72];
  \node[font=\large\bfseries, text=elevelaranja] at (5.18,4.30) {$120^\circ$};
  \node[font=\large\bfseries, text=elevelaranja] at (2.82,4.30) {$120^\circ$};
  \node[font=\large\bfseries, text=elevelaranja] at (4.00,2.20) {$120^\circ$};
  \node[font=\Large\bfseries, text=elevecinza] at (4.00,7.18)
    {$120^\circ+120^\circ+120^\circ=360^\circ$};
\end{tikzpicture}

\end{document}
